\documentclass[conference,a4paper,flushend]{cs-techrep} % (based on IEEEtran.cls)

% The class iaria.cls loads biblatex/biber with correct IARIA settings
% as well as a set of common packages (times, inputenc[utf8], fontenc[T1],
% graphicx, xcolor, url, orcidlink, hyperref, extdash[shortcuts])

% Packages (self-added, not included in cs-techrep.cls):
\usepackage[shortlabels]{enumitem} % Verbesserte Einstellungsmöglichkeiten von Listen. Nicht bei Beamer anwendbar
% \setlist{itemsep=1ex, topsep=1em}
% \usepackage{graphicx} % Required for inserting images
\usepackage{csquotes} % Korrekte Anführungszeichen setzen. Sprachspezifisch möglich. mit \setquotestyle{style} und \enquote{}.
% \usepackage[american]{babel} % Ändert die Darstellung des Datums in amerikanische Darstellung.
% \usepackage[capitalize, nameinlink, american]{cleveref} % Referenzen besser machen. z.B. mit \cref{}
% \usepackage[backend=biber, style=ieee, alldates=iso]{biblatex} % Zitieren
% \usepackage{array} % Tabellen mit \begin{tabular}
	\usepackage{multirow} % Mehrere Zeilen in eine zusammenfassen. für Spalten einfach \multicolumn benutzen
	% \usepackage{booktabs} % Hübschere Tabellen mit \top/mid/cmid/bottomrule
	\usepackage{minted} % Neuere Bibliothek für Code-Blöcke, mit \begin{listing oder minted}
		\usepackage{amsmath} % American Mathematical Society für mathematische Formeln, z.B. mit \begin{equation}
			\usepackage{amssymb}
			\usepackage{amsfonts}
			
			\addbibresource{references.bib}
			% \addbibresource{xampl.bib} % (xampl file provided by the BibTEX base files)
			% ======================================================================
			% EDIT THESE:
			\cstechrepAuthorListTex{Simon Völkl, Johannes Schieder, Sebastian Rosner and Andre Tien Vu\,\orcidlink{0000-0002-1175-2668}}
			\cstechrepAuthorListBib{Simon Völkl, Johannes Schieder, Sebastian Rosner and Andre Tien Vu} %TODO: correct format?
			% Capitalization: https://capitalizemytitle.com/style/Chicago/
			\cstechrepTitleTex{Minimal Working Example of a cs-techrep Paper} %TODO: replace with actual title
			% IF you need manual linebreaks in the title, then clone the title without linebreaks for BibTeX:
			\cstechrepTitleBib{{\cstechrepTitleTex}}
			\cstechrepDepartment{Department of Electrical Engineering, Media, and Computer Science}
			\cstechrepInstitution{Ostbayerische Technische Hochschule Amberg\-/Weiden}
			\cstechrepAddress{Amberg, Germany}
			\cstechrepSeries{Technical Reports}
			\cstechrepYear{2026}
			\cstechrepMonth{7}
			\cstechrepNumber{CL-\cstechrepYear{}-42}
			\cstechrepLang{english}  % en-US
			% ======================================================================
			% DO NOT DELETE AND DO NOT MODIFY THIS:
			\filecontentsForceExpansion|[] % force command expansion inside a filecontents* environment
			\begin{filecontents*}[overwrite]{selfref.bib}
				@TECHREPORT{selfref,
					author = {|cstechrepAuthorListBib},
					title  = {\cstechrepTitleBib},
					institution = {\cstechrepInstitution, \cstechrepDepartment},
					series = {\cstechrepSeries},
					number = {\cstechrepNumber},
					year   = {|cstechrepYear},
					month  = {|cstechrepMonth},
					langid  = {|cstechrepLang},
				}
			\end{filecontents*}
			
			\begin{document}
				\selectlanguage{\cstechrepLang}
				\maketitle
				\begin{abstract}
					
					Modularbeit in Big Data, Cloud and NoSQL. Das ist der englische Abstract (Zusammenfassung). % TODO: replace with actual abstract
					Nicht benennen: Modularbeit, Personen, Studiengang, ...
					Sondern als Elevator Pitch
					
				\end{abstract}
				\begin{IEEEkeywords}
					Cloud, Big Data, NoSQL, Modularbeit % TODO: replace with actual keywords
				\end{IEEEkeywords}
				
				\section{Introduction}
				Here we go. %TODO: replace with actual introduction
				\cite{robocuprescue2025}
				
				Motivation, zusammenfassender Überblick des Anwendungskonzeptes, fachliches Alleinstellungsmerkmal (USP), Benennung von technischen Schlüssel-Bausteinen, Fachvokabular einführen.
				
				\section{Related Work}
				andere Arbeiten oder Technologien (die als Inspiration dienten), Screenshots
				\section{Requirements}
				
				% In die Präambel (ganz oben ins Dokument):
				% \usepackage{enumitem}
				% \usepackage{siunitx}
				
				% -----------------------------------------------------------------------
				
				\subsection*{Project Personas}
				
				To ensure the system meets the diverse needs of all stakeholders, the requirements are guided by the following personas. Each persona has a specific role, defined goals, and distinct interests in the system:
				
				\begin{description}[font=\bfseries, leftmargin=1em]
					\item[Chemist Christian] 
					Focuses on exact environmental conditions. His primary goal is to monitor precise laboratory temperatures to accurately plan experiments and evaluate temperature-dependent reaction rates without interference from climate fluctuations.
					
					\item[Occupational Safety Officer Olivia] 
					Prioritizes employee well-being. She requires immediate, clear indicators of air quality safety to proactively enforce protective measures and prevent workplace health hazards.
					
					\item[Facility Manager Martin] 
					Responsible for building operations and rapid incident response. He needs instant, unmissable alerts regarding critical sensor thresholds to take immediate corrective actions before physical risks escalate.
					
					\item[Project Manager Pia] 
					Looks at the big picture and long-term optimization. She aims to analyze historical data trends to identify recurring patterns, correlate them with facility operations, and implement targeted ventilation strategies.
					
					
				\end{description}
				
				\subsection[US-01: Real-time Temperature Monitoring]{\bfseries US-01: Real-time Temperature Monitoring}
				
				\textit{As Chemist Christian, I want to see the exact current temperature in the laboratory, so that I can accurately plan my experiments and evaluate temperature-dependent reaction rates.}
				
				\smallskip
				\noindent\textbf{Acceptance Criteria:}
				\begin{description}[font=\bfseries]
					\item[AC-01] The dashboard displays the current temperature precisely in \textbf{°C}. Input: latest DB temperature readings. Output: exact numerical temperature value.
					\item[AC-02] The temperature value updates automatically when a new reading arrives — no manual reload required. Test: write new temperature values to DB; display must update within 10\,s.
					\item[AC-03] Displayed temperature values are at most 30\,s old. Test: compare value timestamps with system time; difference $\le$ 30\,s.
					\item[AC-04] If no current temperature data is available (gap $>$ 30\,s), a clearly readable notice \emph{``No current temperature data''} is shown instead of an outdated value.
				\end{description}
				
				\subsection[US-02: Workplace Air Quality Status]{\bfseries US-02: Workplace Air Quality Status}
				
				\textit{As Occupational Safety Officer Oliva, I want to see at a glance whether the air quality in the workspace is currently safe or critical, so that I can immediately initiate protective measures for the employees if necessary.}
				
				\smallskip
				\noindent\textbf{Acceptance Criteria:}
				\begin{description}[font=\bfseries]
					\item[AC-01] The dashboard shows the air quality as a colour-coded indicator based on safety thresholds: \textbf{Green}~(safe), \textbf{Yellow}~(elevated), \textbf{Red}~(critical). Input: latest DB air quality readings. Output: correct status colour.
					\item[AC-02] The status colour updates automatically when a new reading arrives — no manual reload required. Test: write new air quality values to DB; display must update within 5\,s.
					\item[AC-03] The displayed status is based on data that is at most 30\,s old. Test: compare value timestamps with system time; difference $\le$ 30\,s.
					\item[AC-04] If no current air quality data is available (gap $>$ 30\,s), a clearly readable notice \emph{``No current air quality data''} is shown instead of a status colour.
				\end{description}
				
				
				% -----------------------------------------------------------------------
				\subsection[US-03: Critical Alert]{\bfseries US-03: Critical Alert}
				
				\textit{As a facility manager Martin, I want an immediate, unmissable notification the second a sensor reading hits a critical limit, empowering me to take instant corrective action before the situation escalates into a physical health risk.} 
				
				\smallskip
				\noindent\textbf{Acceptance Criteria:}
				\begin{description}[font=\bfseries]
					\item[AC-01] An email notification is sent automatically within 10\,s of a threshold exceedance.
					\item[AC-02] The notification contains: measured value (incl.\ unit), measurement timestamp, and alert type (e.g.\ \emph{``CO critical''}).
					\item[AC-03] Consecutive exceedances of the same type within 5\,min do not trigger a further email (debounce).
					\item[AC-04] The notification is delivered even if the frontend is not open at trigger time.
				\end{description}
				
				
				% -----------------------------------------------------------------------
				\subsection[US-04: 24-Hour History Chart]{\bfseries US-03: 24-Hour History Chart}
				
				\textit{As Project Manager Pia, I want to analyze the sensor data from the past 24 hours in a clear historical chart, so that I can identify recurring patterns of poor air quality, correlate them with specific facility operations, and implement targeted ventilation strategies.}
				
				\smallskip
				\noindent\textbf{Acceptance Criteria:}
				\begin{description}[font=\bfseries]
					\item[AC-01] Dashboard displays a line chart with readings from the last 24\,hours.
					\item[AC-02] Chart is fully interactive within 3\,s of page load.
					\item[AC-03] Periods without data are shown as visible line breaks — no silent interpolation.
					\item[AC-04] Chart updates automatically when new readings arrive, without a page reload.
				\end{description}
				
				
				% -----------------------------------------------------------------------
				\subsection[US-05: Custom History View]{\bfseries US-04: Custom History View}
				
				\textit{As Data Scientist Daniela, I want to extract historical sensor readings with a customizable time range and sampling resolution, so that I can generate precise, high-granularity datasets to train our predictive air-quality models.}
				
				\smallskip
				\noindent\textbf{Acceptance Criteria:}
				\begin{description}[font=\bfseries]
					\item[AC-01] Dashboard provides date pickers for start/end and a resolution selector (hour/day).
					\item[AC-02] A query for 30\,days at daily resolution returns results in under 5\,s.
					\item[AC-03] If an end date before the start date is entered, a clear error message is shown and no query is executed.
					\item[AC-04] Periods without data appear as an explicit \emph{``No data''} row in the table and as a line break in the chart.
				\end{description}
				
				
				% -----------------------------------------------------------------------
				\subsection[US-06: Sensor Status]{\bfseries US-05: Sensor Status}
				
				\textit{As IT Administrator Ian, I want to instantly verify the connection status and the last heartbeat timestamp of each sensor, so that I can proactively detect offline devices and troubleshoot network dropouts before data gaps occur.}
				
				\smallskip
				\noindent\textbf{Acceptance Criteria:}
				\begin{description}[font=\bfseries]
					\item[AC-01] Sensor status (\emph{``Active''}/\emph{``Inactive''}) is prominently visible. Status switches to \emph{``Inactive''} when the last data point is more than 2\,min ago.
					\item[AC-02] Date and time of the last received reading are displayed next to the status (format: \texttt{DD.MM.YYYY HH:MM:SS}).
					\item[AC-03] Status indicator updates automatically — no manual reload required.
					\item[AC-04] If the database connection is temporarily unavailable, the notice \emph{``Connection interrupted''} is shown instead of a false \emph{``Active''} status.
				\end{description}
				
				% -----------------------------------------------------------------------
				\subsection[Definition of Done]{\bfseries Definition of Done}
				% -----------------------------------------------------------------------
				
				A user story is considered complete only when, in addition to its specific acceptance criteria, \emph{all} of the following quality standards are met:
				
				\begin{itemize}[itemsep=2pt, topsep=4pt]
					\item \textbf{Code Review:} Reviewed and approved by at least one other developer (four-eyes principle).
					\item \textbf{Tests:} Unit tests written and passing; coverage $\ge$ 40\,\%.
					\item \textbf{Documentation:} Technical docs and API interfaces updated where applicable.
					\item \textbf{Deployment:} Feature successfully deployed.
					\item \textbf{Functional Test:} Feature manually tested in staging without issues.
				\end{itemize}
				
				
				% -----------------------------------------------------------------------
				\subsection[Prioritisation (MoSCoW)]{\bfseries Prioritisation (MoSCoW)}
				% -----------------------------------------------------------------------
				
				\begin{center}
					\begin{tabular}{@{}ll@{}}
						\toprule
						\textbf{Category}    & \textbf{User Stories} \\
						\midrule
						\textbf{M}ust Have   & US-01 (Current Status Display), US-03 (Critical Alert), US-02 Workplace Air Quality \\[3pt]
						\textbf{S}hould Have & US-06 (Sensor Status) \\[3pt]
						\textbf{C}ould Have  & US-04 (24-Hour History Chart) \\[3pt]
						\textbf{W}on't Have  & US-05 (Custom History View — \textit{deferred to Release~2.0}) \\
						\bottomrule
					\end{tabular}
				\end{center}
				
				
				\section{Minimum Viable Product (MVP): IoT Data Ingestion and Processing via AWS}
				
				\subsection{Objective}
				The MVP implements a server-side IoT pipeline on Amazon Web Services (AWS). It focuses on ingesting simulated sensor data via MQTT, processing and storing it within the AWS ecosystem, and providing access through a Web API. This serves as the technical foundation for the system.
				
				\subsection{System Architecture}
				
				\textbf{Data Generation \& Transmission} \\
				Instead of physical hardware, a Python script simulates realistic sensor data. Messages are transmitted via the MQTT protocol to \textbf{AWS IoT Core}, ensuring a lightweight and software-defined testing environment.
				
				\textbf{Processing Pipeline} \\
				Incoming data is handled by an \textbf{AWS IoT Rule} which triggers a \textbf{Lambda function}. The function validates the payload and persists it in:
				\begin{itemize}
					\item \textbf{Amazon DynamoDB:} Primary storage for structured time-series data.
					\item \textbf{Amazon S3:} Long-term archival of raw data for cost-efficiency.
				\end{itemize}
				
				\textbf{Web API Access} \\
				A REST API, built with \textbf{AWS API Gateway} and \textbf{Lambda}, provides two core endpoints:
				\begin{itemize}
					\item \texttt{GET /live} – Returns the most recent measurement.
					\item \texttt{GET /history} – Retrieves chronological data from DynamoDB.
				\end{itemize}
				
				\subsection{Out of Scope}
				To focus on core infrastructure, the following elements are excluded from this project phase:
				\begin{itemize}
					\item Graphical Frontend/Dashboard for data visualization.
					\item Physical ESP32 hardware and real sensors.
					\item Automated user notifications (e.g., Push or Email alerts).
				\end{itemize}
				% -----------------------------------------------------------------------
				\section{Technical Concept}
				knapp technische Schlüsselmerkmale der Anwendung, rein deskriptiv
				
				
				
				\section{Technical Concept}
				knapp technische Schlüsselmerkmale der Anwendung, rein deskriptiv
				
				% ======== References =========
				\sloppy
				\printbibliography[notcategory=selfref]
				
			\end{document}
